\documentclass[12pt]{article}
\usepackage[T1]{fontenc}
\usepackage[utf8]{inputenc}
\usepackage{lmodern}
\usepackage{textcomp}
\usepackage{enumitem}
\usepackage{geometry}
\usepackage{graphicx}
\usepackage{amsmath, amssymb}
\geometry{margin=1in}
\begin{document}
\begin{center}
Biology - Undergraduate Level Assessment 2 \
Total Marks: 40 \
Answer all questions.
\end{center}

\section*{Multiple Choice (10 marks)}
\begin{enumerate}[label=\arabic*.]
\item In most ecosystems, net primary productivity is important because it represents the
\begin{enumerate}[label=(\alph*)]
    \item biomass of all producers
    \item biomass of all consumers
    \item energy available to heterotrophs
    \item energy consumed by decomposers
    \item total solar energy converted to chemical energy by producers
\end{enumerate}
\item What prevents the stomach from being digested by itsown secretions?
\begin{enumerate}[label=(\alph*)]
    \item The stomach has a high permeability to hydrogen ions
    \item The presence of beneficial bacteria that neutralize stomach acid
    \item The stomach lining is replaced every month
    \item The stomach lining secretes digestive enzymes only when food is present
    \item The production of bile in the stomach neutralizes the digestive acids
\end{enumerate}
\item In modern terminology, diversity is understood to be a result of genetic variation. Sources of variation for evolution include all of the following except
\begin{enumerate}[label=(\alph*)]
    \item mistakes in translation of structural genes.
    \item mistakes in DNA replication.
    \item translocations and mistakes in meiosis.
    \item recombination at fertilization.
\end{enumerate}
\item Charles Darwin was the first person to propose
\begin{enumerate}[label=(\alph*)]
    \item that humans evolved from apes.
    \item that evolution occurs.
    \item that evolution occurs through a process of gradual change.
    \item that the Earth is older than a few thousand years.
    \item a mechanism for how evolution occurs.
\end{enumerate}
\item How can you determine whether or not a given population is in genetic equilibrium?
\begin{enumerate}[label=(\alph*)]
    \item Compare gene frequencies and genotype frequencies of two generations.
    \item Evaluate the population's response to environmental stress
    \item Assess the mutation rates of the population's genes
    \item Compare the population size over generations
    \item Calculate the ratio of males to females in the population
\end{enumerate}
\end{enumerate}

\section*{True / False (10 marks)}
\textit{Answer True or False}
\begin{enumerate}[label=\arabic*.]
\setcounter{enumi}{5}
\item True or False: If messenger RNA molecules were not destroyed, they would replicate rapidly, leading to excessive protein synthesis.
\item True or False: Autoradiography is a technique used to visualize amino acids incorporated into proteins in cells.
\item True or False: Normal somatic cells of a human male contain two active X chromosomes.
\item True or False: Viruses can reproduce independently without the need for a host cell.
\item True or False: The symplastic pathway for sucrose movement involves the bundle sheath, phloem parenchyma, companion cell, and sieve tube.
\end{enumerate}

\section*{Long Form Response (20 marks)}
\textit{Answer each question with a clear and complete explanation.}
\begin{enumerate}[label=\arabic*.]
\setcounter{enumi}{10}
\item Discuss the emergence of the first mammals during the Carboniferous Period and analyze the key differences between monotremes and marsupials in terms of their reproductive strategies and adaptations.
\item Discuss the challenges associated with classifying tissues in higher plants, considering their evolutionary heritage and the implications for understanding plant biology.
\end{enumerate}

\vfill
\noindent\textit{End of Paper}
\end{document}